\begin{definition}[Finite signed permutation representations]%
    \label{def:peps_signed_permutation_representation}
    \lean{TNLean.signedPermutationSubgroup,
        TNLean.instFiniteSignedPermutationSubgroup,
        TNLean.signedPermutationRepresentation,
        TNLean.signedPermutationRepresentation_mem_unitaryGroup,
        TNLean.signedPermutationRepresentation_range}
    \leanok
    For a finite index set $I$, let $G_I$ be the group of matrices sending
    each standard basis vector to another standard basis vector with sign
    $+1$ or $-1$. Its defining representation is unitary, and $G_I$ is finite.
    These groups supply the finite-dimensional matrix representations used
    in the construction of \cite[Theorem~4.1]{Schuch2010PEPS}.
\end{definition}

\begin{theorem}[Matrix units from signed permutations]%
    \label{thm:peps_signed_permutation_span}
    \lean{TNLean.exists_signedPermutationRepresentation_sub_eq_single,
        TNLean.span_signedPermutationRepresentation}
    \leanok
    \uses{def:peps_signed_permutation_representation}
    Every matrix unit $E_{ij}$ is of the form $(P-Q)/2$ for two signed
    permutation matrices $P,Q\in G_I$. Consequently,
    $\operatorname{span}_{\mathbb C}G_I=M_I(\mathbb C)$.
\end{theorem}

\begin{proof}\leanok
    Choose a permutation taking $j$ to $i$, and let $P$ be its permutation
    matrix. Negate only its column $j$ to obtain $Q$. Then
    $P-Q=2E_{ij}$. The matrix units span all matrices. The empty index set
    causes no exception: its matrix space is zero.
\end{proof}

\begin{definition}[Independent groups on multiplicity factors]%
    \label{def:peps_block_signed_permutation_representation}
    \lean{TNLean.blockSignedPermutationRepresentation,
        TNLean.blockSignedPermutationRepresentation_mem_unitaryGroup}
    \leanok
    \uses{def:peps_signed_permutation_representation}
    For finitely many pairs of dimensions $(m_k,d_k)$, the finite group
    $G=\prod_kG_{m_k}$ has the unitary representation
    $U(g)=\bigoplus_k g_k\otimes I_{d_k}$ on
    $\bigoplus_k\mathbb C^{m_k}\otimes\mathbb C^{d_k}$.
\end{definition}

\begin{theorem}[Commutant of independent multiplicity actions]%
    \label{thm:peps_block_signed_permutation_commutant}
    \lean{TNLean.commutes_blockSignedPermutationRepresentation_iff}
    \leanok
    \uses{def:peps_block_signed_permutation_representation,
        thm:peps_signed_permutation_span}
    A matrix commutes with every $U(g)$ if and only if it is of the form
    $\bigoplus_k I_{m_k}\otimes B_k$, with $B_k\in M_{d_k}(\mathbb C)$.
\end{theorem}

\begin{proof}\leanok
    Two group elements which agree on all factors except $k$ produce,
    by their half-difference, $E_{ij}\otimes I_{d_k}$ supported on block $k$.
    Thus the span of $U(G)$ contains every
    $\bigoplus_k C_k\otimes I_{d_k}$. Commutation passes to this span.
    Swap the two tensor factors on each block and apply the commutant
    formula for a direct sum of full matrix algebras. Conversely,
    $I_{m_k}\otimes B_k$ commutes with $g_k\otimes I_{d_k}$.
\end{proof}

\begin{theorem}[Unital matrix algebras as finite-group commutants]%
    \label{thm:peps_finite_group_commutant_unital}
    \lean{StarSubalgebra.exists_finite_unitary_group_commutant}
    \leanok
    \uses{thm:peps_block_signed_permutation_commutant}
    Let $S\subset M_I(\mathbb C)$ be a unital C*-subalgebra, with $I$ finite.
    There exist a finite group $G$ and a unitary representation $U$ such that
    $S=\{A\in M_I(\mathbb C): AU(g)=U(g)A\text{ for every }g\in G\}$.
    This is the unital version of \cite[Theorem~4.1]{Schuch2010PEPS}.
    The source's definition of C*-algebra does not require the identity;
    the printed nonunital assertion fails because every commutant contains
    the identity. See \cite{gap:scp10_finite_group_commutant_unital}.
\end{theorem}

\begin{proof}\leanok
    The unitary block-decomposition theorem for matrix C*-algebras
    \cite[Theorem~6.14]{Wolf2012Quantum} gives a unitary $W$ with
    $W^*SW=\bigoplus_k I_{m_k}\otimes M_{d_k}(\mathbb C)$.
    Apply the independent signed permutation groups to the multiplicity
    factors. Their commutant is this block algebra. Conjugating their
    representation by $W$ therefore gives a finite unitary group with
    commutant $S$. All decomposition data are obtained from $S$.
\end{proof}


\begin{theorem}[C*-algebra tensors have a finite group symmetry]%
    \label{thm:peps_cstar_span_g_injective_mps}
    \lean{StarSubalgebra.exists_isGInjective_mpsSiteMap_of_span}
    \leanok
    \uses{thm:peps_finite_group_commutant_unital,
        def:peps_mps_g_injective}
    Suppose that a family of complex matrices $A^i$ spans a unital
    C*-subalgebra $S\subset M_I(\mathbb C)$, with $I$ finite.
    There exist a finite group $G$ and a unitary representation $U$ whose
    commutant is $S$, such that the tensor $A$ is $G$-injective.
    This gives the passage from the unital version of
    \cite[Theorem~4.1]{Schuch2010PEPS} to its Definition~4.2.
    The unital restriction is recorded in
    \cite{gap:scp10_finite_group_commutant_unital}.
\end{theorem}

\begin{proof}\leanok
    Choose the finite unitary representation with commutant $S$.
    Each $A^i\in S$ commutes with this representation, so its site map
    $X\mapsto (\operatorname{tr}(A^iX))_i$ is invariant under conjugation.
    If $X\in S$ is in the kernel, linearity gives
    $\operatorname{tr}(YX)=0$ for every $Y\in S$.
    Since $S$ is closed under adjoints, take $Y=X^*$.
    Then $\operatorname{tr}(X^*X)=0$, hence $X=0$.
    Thus the site map is injective on the invariant subspace $S$.
\end{proof}
